\section{Measure and integration from stable finite observations}
\label{sec:measure}

The measure-theoretic improvement is to use two compatible distances on the
same finite data.  The strong distance controls integration; a bounded
distance controls convergence in measure.  A quantitative tail estimate
converts the second into the first.  This produces limits rather than
assuming an integrable limit in advance.  A second construction recovers a
density from dyadic masses.  Its convergence certificate can be derived from
local oscillation bounds on those masses, so a density evaluator is not an
input to the construction.

\subsection{The rational dyadic interface}

\begin{definition}[Dyadic step data]\label{m:step}
A step function of level \(r\) is a vector
\(a=(a_0,\ldots,a_{2^r-1})\) of rational numbers.  Its intended value is
\(a_i\) on \([i2^{-r},(i+1)2^{-r})\).  The right endpoint of the unit interval
is immaterial.  Refinement from level \(r\) to level \(s\ge r\) repeats each
entry \(2^{s-r}\) times.  Define
\[
 I(a)=2^{-r}\sum_{i<2^r}a_i,\qquad
 \|a\|_1=2^{-r}\sum_{i<2^r}|a_i|,\qquad
 Q(a)=2^{-r}\sum_{i<2^r}a_i^2.
\]
For two steps, first refine to the larger level.  All three quantities are
unchanged by refinement.
\end{definition}

This definition is an arithmetic data type.  Signed integers encode rational
numerators, and positive natural numbers encode denominators.  Comparison,
refinement, finite summation, and the functions
\(\min,\max,|\cdot|\) have primitive recursive implementations.  In
particular, the length \(2^r\) is an explicit loop bound.  None of these
operations requires selecting a point of an abstract measurable space.

\begin{lemma}[Finite integral estimates]\label{m:finite}
For rational steps on a common refinement,
\[
 |I(a)|\le\|a\|_1,\qquad
 \|a+b\|_1\le\|a\|_1+\|b\|_1,\qquad
 \bigl|\|a\|_1-\|b\|_1\bigr|\le\|a-b\|_1,
 \qquad
 \|a\|_1^2\le Q(a).
\]
The integral is rational linear and preserves entrywise order.
\end{lemma}

\begin{proof}
The first three assertions follow by summing the corresponding rational
inequalities entry by entry.  For the last, write \(J=2^r\) and \(u_i=|a_i|\).
The exact identity
\[
 J\sum_{i<J}u_i^2-\Bigl(\sum_{i<J}u_i\Bigr)^2
   =\sum_{0\le i<j<J}(u_i-u_j)^2
\]
has a nonnegative right side.  Divide by \(J^2>0\).
Linearity and monotonicity are finite distributivity and addition of
inequalities.  Alternatively, clear the common denominator throughout; the
proof is then an integer polynomial identity and natural-number inequalities.
\end{proof}

\begin{definition}[Strong names]\label{m:l1}
An \(L^1\) name is a fixed primitive recursive producer \(p\mapsto a_p\) of
dyadic steps, together with proofs of
\[
 q\ge p\quad\Longrightarrow\quad \|a_q-a_p\|_1\le2^{-p}.
 \tag{M1}
\]
Equivalent names have \(L^1\) distance zero, witnessed by explicit precision
maps.  For example, a certificate of equivalence may supply a primitive
recursive \(e(k)\) and prove
\(\|a_p-b_q\|_1\le2^{-k}\) whenever \(p,q\ge e(k)\).
We use \(L^1_{\rm dyad}\) for this represented completion.
\end{definition}

Normalization (M1) is slightly stronger than the core regular-name bound
\(\|a_p-a_q\|\le2^{-p}+2^{-q}\).  Replacing \(a_p\) by \(a_{p+1}\) converts
that bound into (M1) for \(q\ge p\); the representations therefore differ
only by a fixed precision shift.

The notation \(f\in L^1_{\rm dyad}\) describes a represented object.  In the
object theory, each theorem is about the fixed producers, rational
inequalities, and their certificates.  It is not a quantification over a
new set of all functions.  The finite estimates give names for
\(I(f)\) and \(\|f\|_1\).  The notation
\(\|f-a_p\|_1\le2^{-p}\) abbreviates the corresponding regular-real
inequality obtained from (M1).  Linearity, positivity, and the integral
estimate extend by applying the finite results with guard bits.
Lattice operations such as absolute value and positive part extend by
their Lipschitz estimates.  In particular, the order \(u\le v\) means
\(\|(u-v)^+\|_1=0\).  This defines the domination hypothesis used below
without selecting pointwise representatives.

\subsection{Boolean completion and effective countable additivity}

A dyadic set is a step vector whose entries belong to \(\{0,1\}\).
Its measure is the average of its entries, and
\[
 d(A,B)=I(|1_A-1_B|)=\lambda(A\mathbin{\triangle}B).
\]
Complement is an isometry.  Union and intersection satisfy
\[
 d(A\cup B,A'\cup B')\le d(A,A')+d(B,B'),
 \qquad
 d(A\cap B,A'\cap B')\le d(A,A')+d(B,B').
 \tag{M2}
\]
These follow by checking the finitely many Boolean possibilities at one
entry and summing.  Consequently they act on strong Boolean names.
This is an effective measure algebra: equality means zero measure of the
symmetric difference and does not decide pointwise membership.

\begin{theorem}[Effective unions]\label{m:union}
Let \(A_i\) be a uniformly presented sequence of Boolean names.  Suppose a
primitive recursive \(T\) and finite certificates give
\[
 m\ge n\ge T(p)
 \quad\Longrightarrow\quad
 \lambda\Bigl(\bigcup_{n\le i<m} A_i\Bigr)\le2^{-p}.
 \tag{M3}
\]
Then the union of the sequence has a Boolean name.  If the \(A_i\) are
pairwise disjoint modulo zero measure, its measure is the regular-real sum
\[
 \lambda\Bigl(\bigcup_{i=0}^{\infty}A_i\Bigr)
       =\sum_{i=0}^{\infty}\lambda(A_i).
 \tag{M4}
\]
A sufficient input for (M3) is the stronger bound
\(\sum_{n\le i<m}\lambda(A_i)\le2^{-p}\).
\end{theorem}

\begin{proof}
Put \(U_n=\bigcup_{i<n}A_i\).  Boolean algebra gives
\(U_m\mathbin{\triangle}U_n\subseteq\bigcup_{n\le i<m}A_i\) for \(m\ge n\).
Thus (M3) is exactly a Cauchy estimate for the prefix unions.

For an explicit finite output at precision \(p\), let
\(N=T(p+2)\) and approximate each \(A_i\), \(i<N\), with error at most
\(2^{-h}\), where
\[
 h=p+2+\lceil\log_2(\max(1,N))\rceil.
\]
Take the finite union of those approximants.  Repeated application of (M2)
bounds its error from \(U_N\) by \(N2^{-h}\le2^{-p-2}\).
For two requested precisions \(p,q\), compare both finite outputs with the
finite prefix \(U_{\max(T(p+2),T(q+2))}\).  Condition (M3) bounds the two
prefix discrepancies by \(2^{-p-2}\) and \(2^{-q-2}\), respectively.
The two approximation errors and these two tail errors show that outputs
at \(p\le q\) differ by at most
\(2^{-p-1}+2^{-q-1}\le2^{-p}\).  This constructs the strong Boolean name
without first assuming that an infinite union exists.
The case \(N=0\) is the empty union.

For disjoint sets, the finite identity
\(\lambda(U_n)=\sum_{i<n}\lambda(A_i)\) follows from Boolean additivity.
Apply the same tail estimate to both sides.  This proves (M4) with an actual
regular-real modulus, rather than an unmodulated supremum.
The word union also has its order-theoretic meaning.  Define \(A\le B\) by
\(\lambda(A\setminus B)=0\).  Every \(A_i\) lies below each sufficiently long
prefix and hence below the constructed limit.  Any upper bound of all
\(A_i\) bounds every finite prefix.  Continuity of Boolean difference and
measure then makes it an upper bound of the limit.  The constructed union
is therefore the least upper bound in the completed Boolean algebra.
\end{proof}

The adjective \emph{uniformly} above includes a bound for producing the
variable-length list of approximants.  For fixed primitive recursive
producers, the bounded loop is primitive recursive.  A relational
presentation instead supplies one batch trace and its verification
certificate.  Individual totality assertions for each \(A_i\) alone do not
license this batch construction.

\subsection{Fubini and conditional averages}

On the square, a rational rectangle step is an array
\(a_{ij}\), \(i<2^r,\ j<2^s\), with integral and \(L^1\) norm given by the
average of its entries and absolute entries.  Define the partial average
\[
 (P_y a)_i=2^{-s}\sum_{j<2^s}a_{ij}.
\]

\begin{theorem}[Fubini as an operator on names]\label{m:fubini}
Partial averaging is an \(L^1\) contraction:
\[
 \|P_y a-P_y b\|_1\le\|a-b\|_1.
 \tag{M5}
\]
It therefore sends each strong rectangle name directly to a strong
one-dimensional name.  On these completions,
\[
 I_x(P_y f)=I_{x,y}(f)=I_y(P_x f).
 \tag{M6}
\]
\end{theorem}

\begin{proof}
For each row, the absolute value of an average is no larger than the average
of absolute values.  Average this row inequality over \(i\) to obtain (M5).
For a finite array, both iterated integrals equal
\(2^{-r-s}\sum_{i,j}a_{ij}\).  Apply that exact identity to the \(p\)-th
array in a strong name.  By (M5), all three scalar names have the same
finite values and compatible moduli.  This proves (M6).
\end{proof}

This formulation constructs the iterated integral without selecting
pointwise sections on null exceptional sets.  A separate comparison with a
classical measure-space interpretation may identify it with the usual
Fubini integral.  No claim that every individual section is integrable is
needed internally.

For \(n\in\mathbb N\), define \(E_n a\) to be the average of \(a\) on each
level-\(n\) dyadic cell.  If \(a\) has level below \(n\), simply refine it.
If it has level \(r>n\), average each block of \(2^{r-n}\) entries.

\begin{theorem}[Conditional averages with a derived rate]\label{m:tower}
The operators \(E_n\) preserve the integral and positivity, are \(L^1\)
contractions, and satisfy
\[
 E_nE_m=E_{\min(n,m)}.
 \tag{M7}
\]
Every strong name \(f=(a_p)\) has a primitive recursive convergence modulus
for \(E_n f\to f\).  In fact, if \(\ell(a)\) denotes the level of a step,
\[
 n\ge\ell(a_{p+1})
 \quad\Longrightarrow\quad
 \|E_n f-f\|_1\le2^{-p}.
 \tag{M8}
\]
\end{theorem}

\begin{proof}
Each property in (M7) is a finite identity of block averages.  In particular,
averaging twice over nested dyadic partitions averages over the coarser
partition.  The contraction follows exactly as in (M5).  If
\(n\ge\ell(a_{p+1})\), then \(E_n a_{p+1}=a_{p+1}\), and hence
\[
 \|E_nf-f\|_1
 \le\|E_n(f-a_{p+1})\|_1+\|a_{p+1}-f\|_1
 \le2\|f-a_{p+1}\|_1\le2^{-p}.
 \]
The displayed level is computed from one finite approximation.  Thus no
independent martingale-convergence oracle is required once a strong
\(L^1\) name is already available.
\end{proof}

\subsection{Stable tests, clipping, and the repaired Vitali estimate}

Define, for rational \(x,y\) and \(K\ge1\),
\[
 c_K(x)=\max(-K,\min(K,x)),\qquad
 t_K(x)=\max(|x|-K,0),\qquad
 \rho(x,y)=\min(1,|x-y|).
\]
On finite steps put
\[
 d_0(a,b)=I(\rho(a,b)),\qquad \tau_K(a)=I(t_K(a)).
\]
The functions \(t_K\) and \(x\mapsto c_K(x)\) are \(1\)-Lipschitz, and
\(\rho\) is \(1\)-Lipschitz in each argument.  These quantities consequently
extend to \(L^1\) names with explicit moduli.  Their use avoids deciding a
sharp threshold event for an arbitrary real-valued name.  For finite steps,
\[
 \lambda(|a-b|>\eta)\le \frac{d_0(a,b)}{\min(1,\eta)}
 \quad(\eta>0),
 \qquad
 d_0(a,b)\le\eta+\lambda(|a-b|>\eta)\quad(0<\eta\le1).
 \tag{M9}
\]
The first inequality follows entry by entry on the indicated event; the
second splits that event from its complement.  Thus \(d_0\) is a stable
quantitative substitute for convergence in measure.
The inequality
\(\min(1,a+b)\le\min(1,a)+\min(1,b)\) for \(a,b\ge0\) proves the triangle
inequality for \(d_0\).  Moreover \(d_0\le\|\cdot\|_1\).
On \(L^1\) names it separates equivalence classes: the excess tails of each
fixed name can be made arbitrarily small by Proposition~\ref{m:ui}, and
(M10) then turns \(d_0(f,g)=0\) into \(\|f-g\|_1=0\).

\begin{lemma}[Clipping kernel]\label{m:clipping}
For rational steps and \(K\ge1\),
\[
 \|a-b\|_1
 \le 2K\,d_0(a,b)+\tau_K(a)+\tau_K(b).
 \tag{M10}
\]
The same inequality holds for strong \(L^1\) names.
\end{lemma}

\begin{proof}
First, \(|x-c_K(x)|=t_K(x)\).  The triangle inequality gives
\[
 |x-y|\le t_K(x)+|c_K(x)-c_K(y)|+t_K(y).
 \]
Clipping is nonexpansive and its image has diameter \(2K\), so its middle
term is at most \(\min(|x-y|,2K)\).  If \(|x-y|\le1\), this is at most
\(2K|x-y|\), since \(2K\ge1\).  If \(|x-y|\ge1\), it is at most \(2K\).
In both cases the bound is \(2K\rho(x,y)\).  Sum the resulting pointwise
inequality.  The passage to names uses only the already proved Lipschitz
estimates, so it preserves the rational bound without threshold-event
selection.
\end{proof}

\paragraph{A precise repair of a repository argument.}
The \emph{Vitali form} proof in the inspected
\texttt{expected\_value\_ipra\_2026-11-16} packet, section
\texttt{12\_convergence.tex}, displays the shortcut
\[
 |x-y|\ \le\
 |x-y|1_{\{|x|\le K,\ |y|\le K\}}
 +|x|1_{\{|x|>K\}}+|y|1_{\{|y|>K\}}.
 \tag{M11}
\]
This displayed inequality is false: \(x=2,\ y=-1,\ K=1\) gives \(3\le2\).
The omitted contribution is the bounded coordinate on a region where only
the other coordinate exceeds the threshold.  Clipping gives a correct proof
of the intended conclusion, with the excess tails \(t_K\).  This correction
does not invalidate the repository's separate common-envelope theorem
\texttt{Probability/TruncationEstimate.lean}.  That module proves the
different, valid estimate with \(|X-Y|\le2G\) and a tail of the common
envelope \(G\); its own documentation explicitly leaves normalization and
refinement outside its derived result.

\subsection{A limit-producing quantitative Vitali theorem}

\begin{theorem}[Vitali without a pre-supplied limit]\label{m:vitali_cauchy}
Let \(f_n\) be a fixed primitive recursive sequence of rational dyadic
steps.  Suppose there are primitive recursive functions \(N,U\), with
\(U(p)\ge1\), and proofs of
\[
 m,n\ge N(q)\ \Longrightarrow\ d_0(f_m,f_n)\le2^{-q},
 \qquad
 \tau_{U(p)}(f_n)\le2^{-p}\quad\hbox{for every }n.
 \tag{M12}
\]
Define
\[
 K_p=U(p+2),\quad
 h_p=\lceil\log_2 K_p\rceil,\quad
 M(p)=N(p+3+h_p),\quad
 \widehat M(p)=\max_{j\le p}M(j).
 \tag{M13}
\]
The logarithm notation means bounded search for the least integer \(h\)
with \(K_p\le2^h\).  Then
\[
 m,n\ge M(p)\quad\Longrightarrow\quad
 \|f_m-f_n\|_1\le\tfrac34\,2^{-p}.
 \tag{M14}
\]
In particular \(a_p=f_{\widehat M(p)}\) is an \(L^1\) name, and
\(f_n\) converges to its represented limit with an explicit \(L^1\) modulus.
\end{theorem}

\begin{proof}
Use (M10) with \(K=K_p\).  The two excess tails are each at most
\(2^{-p-2}\).  For \(m,n\ge M(p)\), the metric contribution is at most
\[
 2K_p\,2^{-p-3-h_p}\le2^{-p-2}.
 \]
This proves (M14).  If \(q\ge p\), then both
\(\widehat M(p)\) and \(\widehat M(q)\) are at least \(M(p)\), proving
\(\|a_q-a_p\|_1\le(3/4)2^{-p}\le2^{-p}\).
The limiting object has therefore been constructed by a concrete producer.

To prove convergence of the original sequence, fix \(n\ge M(p+1)\).
For every \(q\ge p+1\), (M14) bounds
\(\|f_n-a_q\|_1\) by \((3/4)2^{-p-1}\).
The resulting regular-real inequality is smaller than \(2^{-p}\).
Since \(|I(f_n)-I(a_q)|\le\|f_n-a_q\|_1\), the same schedule controls
convergence of expectations.  The proof uses no choice of a subsequence
converging pointwise and no Fatou lemma.
\end{proof}

The theorem also applies to uniformly presented \(L^1\) names.  First use
(M14) for those names, then let the output step be an approximation to
\(f_{\widehat M(p+2)}\) with \(L^1\) error at most \(2^{-p-3}\).
For two output precisions \(q\ge p\), the two approximation errors and
(M14) sum to at most
\[
 2^{-p-3}+2^{-q-3}+\tfrac34\,2^{-p-2}<2^{-p}.
 \]
The batch-production requirement remains explicit.

\begin{proposition}[Two producers for uniform integrability]\label{m:ui}
\begin{enumerate}
\item If \(Q(f_n)\le C\) for every \(n\), then
\[
 \tau_K(f_n)\le C/(4K)\qquad(K>0).
 \tag{M15}
\]
For any integer \(C_0\ge\max(1,C)\),
\(U(p)=C_0\,2^p\) satisfies the tail part of (M12).
\item If \(|f_n|\le G\) and a strong \(L^1\) name \(G=(g_p)\) is supplied,
the latter name computes a valid \(U\).  It suffices to take
\[
 U(p)=\max\bigl(1,\max_i |(g_{p+1})_i|\bigr).
 \tag{M16}
\]
Thus a strong dominating name supplies its own amplitude-tail modulus.
\end{enumerate}
\end{proposition}

\begin{proof}
For \(u\ge0\), the inequality \(4K(u-K)^+\le u^2\) is immediate when
\(u\le K\); otherwise it is equivalent to \((u-2K)^2\ge0\).
Average it to prove (M15).
For the second claim, put \(B=\max_i|(g_{p+1})_i|\).  Since \(|g_{p+1}|\le B\),
\[
 (|G|-B)^+\le |G-g_{p+1}|,
 \qquad
 \tau_{U(p)}(G)\le2^{-p-1}.
 \]
The map \(t_K\) is increasing in \(|x|\), so domination gives the same bound
for every \(f_n\).  The argument is a Lipschitz estimate on names and does
not require evaluating \(G\) pointwise.
More explicitly, the finite inequality
\[
 t_K(x)\le(|x|-y)^+ +(y-K)^+
          \le(|x|-y)^+ +t_K(y)
 \]
passes to names.  The order definition
\(\|(|f_n|-G)^+\|_1=0\) removes the first term and yields
\(\tau_K(f_n)\le\tau_K(G)\).
\end{proof}

\subsection{Density construction from masses}

\begin{definition}[Coherent dyadic mass data]\label{m:masses}
Supply nonnegative rational numbers \(v_{n,i}\), \(i<2^n\), produced by a
fixed primitive recursive algorithm, with exact coherence
\[
 v_{n,i}=v_{n+1,2i}+v_{n+1,2i+1}.
 \tag{M17}
\]
Their dyadic density steps and energies are
\[
 d_{n,i}=2^n v_{n,i},\qquad H_n=Q(d_n).
 \]
Write \(\nu(A)\) for the finite sum of masses of any dyadic set \(A\) at a
common level.  Coherence makes this independent of refinement.
\end{definition}

\begin{lemma}[Orthogonal energy identity]\label{m:energy}
For \(m\ge n\),
\[
 E_n d_m=d_n,\qquad
 Q(d_m-d_n)=H_m-H_n\ge0,\qquad
 \|d_m-d_n\|_1^2\le H_m-H_n.
 \tag{M18}
\]
\end{lemma}

\begin{proof}
Repeated use of (M17) says that the average of the finer densities over each
coarse cell is its coarse density.  This is the first assertion.
Consequently
\[
 I(d_m d_n)=I\bigl((E_nd_m)d_n\bigr)=I(d_n^2)=H_n.
 \]
Expand \(I((d_m-d_n)^2)\) and cancel the cross term to obtain the second
assertion.  Nonnegativity follows because a square has nonnegative entries.
The last assertion is the finite Cauchy--Schwarz estimate of
Lemma~\ref{m:finite}.  All products in this proof are finite products of
rational steps; no unrestricted multiplication map on \(L^1\) is asserted.
\end{proof}

\begin{theorem}[Effective Radon--Nikodym fragment]\label{m:rn}
Suppose coherent dyadic masses have an energy-tail certificate
\[
 m\ge n\ge R(p)\quad\Longrightarrow\quad H_m-H_n\le4^{-p}
 \tag{M19}
\]
for a primitive recursive \(R\).  Let
\[
 r(p)=\max\bigl(p,\max_{j\le p}R(j)\bigr).
 \]
Then \(p\mapsto d_{r(p)}\) is a nonnegative \(L^1\) name \(f\), and
\[
 \nu(A)=\int_A f\,d\lambda
 \tag{M20}
\]
for every dyadic set \(A\).  The represented density is unique among
\(L^1_{\rm dyad}\) names satisfying (M20).
\end{theorem}

\begin{proof}
If \(q\ge p\), then \(r(q)\ge r(p)\ge R(p)\).  By (M18),
\(\|d_{r(q)}-d_{r(p)}\|_1^2\le4^{-p}\).
Both sides are nonnegative rationals, so
\(\|d_{r(q)}-d_{r(p)}\|_1\le2^{-p}\); a proof by multiplication of
inequalities avoids a square-root operation.  Each approximant is
nonnegative.

If \(A\) has level \(k\), then
\(\int_A d_n\,d\lambda=\nu(A)\) whenever \(n\ge k\), by (M17).
The schedule \(r(p)\ge p\) eventually exceeds \(k\).
Moreover, integration over \(A\) is an \(L^1\) contraction.
Thus the scalar name for the left integral in (M20) is the constant name
\(\nu(A)\).

If \(f\) and \(g\) both represent the masses, then the cell integrals show
\(E_n f=E_n g\) for every \(n\).  Theorem~\ref{m:tower} supplies convergence
rates for both conditional averages.  At a precision where both errors are
at most \(2^{-p-1}\),
\[
 \|f-g\|_1\le\|f-E_nf\|_1+\|E_ng-g\|_1\le2^{-p}.
 \]
This is effective equality of the density names.  No Lebesgue differentiation
theorem is used.
\end{proof}

Condition (M19) already gives a global energy bound:
\(C=H_{R(0)}+1\) bounds \(H_n\), by monotonicity for \(n<R(0)\) and
the precision-zero tail estimate for \(n\ge R(0)\).
Choose an integer \(J\ge1\) with \(J^2\ge C\).
Finite Cauchy--Schwarz gives
\[
 |\nu(A)-\nu(B)|
 \le J\sqrt{\lambda(A\mathbin{\triangle}B)}.
 \tag{M21}
\]
The square root is only notation for the squared rational inequality.
It supplies a modulus that extends \(\nu\) to the effective measure algebra:
to obtain value error \(2^{-p}\), request set error at most
\(4^{-p}/J^2\).  The density integral extends with the same modulus, and
(M20) holds there as well.  Explicitly, if \(J\le2^h\) and \(A_p\) is a
rapid Boolean name, then \(p\mapsto\nu(A_{2p+2h})\) is a rapid scalar name.
The integral is first defined on dyadic sets by its \(L^1\) contraction and
then extended by (M21); this does not introduce unrestricted multiplication
of \(L^1\) names.  Countable additivity is obtained for the
tail-controlled unions of Theorem~\ref{m:union}, with the measure tail
converted by (M21): a \(\lambda\)-tail cutoff \(T\) yields the
\(\nu\)-tail cutoff \(p\mapsto T(2p+2h)\).
This is a represented Radon--Nikodym theorem on the
specified effective algebra, not an assertion about arbitrary
\(\sigma\)-algebras.

\begin{theorem}[A local producer for the energy modulus]\label{m:oscillation}
An energy-tail certificate need not be an oracle.  Suppose the dyadic
densities satisfy the uniform sibling bound
\[
 |d_{n+1,2i}-d_{n+1,2i+1}|\le L\,2^{-n}
 \quad(i<2^n)
 \tag{M22}
\]
for a supplied nonnegative rational \(L\).  Then
\[
 H_m-H_n\le\frac{L^2}{3}\,4^{-n}\quad(m\ge n).
 \tag{M23}
\]
Choose \(h\) with \(L\le2^h\).  The primitive recursive schedule
\(R(p)=p+h\) satisfies (M19), so the masses construct their own density name.
\end{theorem}

\begin{proof}
For siblings \(u,v\), their parent density is \((u+v)/2\), and
\[
 \frac{u^2+v^2}{2}-\left(\frac{u+v}{2}\right)^2
   =\frac{(u-v)^2}{4}.
 \]
Average this identity over the level-\(n\) cells and apply (M22):
\(H_{n+1}-H_n\le(L^2/4)4^{-n}\).
For \(m\ge n\), telescope this inequality.  The finite geometric sum gives
\[
 H_m-H_n
 \le\frac{L^2}{4}\sum_{j=n}^{m-1}4^{-j}
 =\frac{L^2}{3}4^{-n}(1-4^{-(m-n)})
 \le\frac{L^2}{3}4^{-n}.
 \]
At \(n\ge p+h\), the right side is at most \(4^{-p}/3\), hence at most
\(4^{-p}\).  The same calculation also derives the global energy bound
\(H_n\le H_0+L^2/3\).
\end{proof}

\paragraph{An explicit example from masses.}
For a dyadic interval \([a,b)\), supply
\[
 \nu([a,b))=\frac23(b-a)+\frac13(b^2-a^2).
 \]
These are rational, positive, coherent masses of total one.  Direct finite
calculation gives
\[
 d_{n,i}=\frac23+\frac{2i+1}{3\cdot2^n},
 \qquad
 H_n=\frac{28}{27}-\frac{4^{-n}}{27},
 \qquad
 |d_{n+1,2i+1}-d_{n+1,2i}|=\frac13\,2^{-n}.
 \]
Thus the construction produces the density \(f(t)=2(1+t)/3\), with
\(H_\infty-H_n=4^{-n}/27\).  The formula for \(f\) verifies the interpretation
after the construction; the algorithm itself needs only the masses and the
sibling certificate.

\subsection{What an energy bound alone does give}

\begin{theorem}[Finite energy metastability]\label{m:metastability}
Assume only \(H_n\le C_0\), where \(C_0\) is a positive integer.
For every fixed primitive recursive challenge \(g\) and precision \(p\),
one can primitive recursively find \(N\) such that
\[
 m,n\in[N,\max(N,g(N))]
 \quad\Longrightarrow\quad
 \|d_m-d_n\|_1\le2^{-p}.
 \tag{M24}
\]
An explicit search bound is obtained by setting
\[
 t_0=0,\quad t_{j+1}=\max(t_j,g(t_j)),\quad
 B=C_0\,4^p+1
 \]
and searching \(j<B\) for
\(H_{t_{j+1}}-H_{t_j}\le4^{-p}\), then taking \(N=t_j\).
\end{theorem}

\begin{proof}
If all \(B\) increments were larger than \(4^{-p}\), telescoping would give
\[
 H_{t_B}-H_0>B\,4^{-p}>C_0,
 \]
contrary to \(0\le H_0\le H_{t_B}\le C_0\).
The search therefore succeeds.  Monotonicity of \(H\) implies that any two
levels in the selected interval have energy difference at most the endpoint
increment.  Apply (M18).
The recursion for \(t_j\) is primitive recursion using the fixed function
symbol for \(g\).  No evaluator uniform over all primitive recursive program
codes has been introduced.
\end{proof}

This is a substantive finite conclusion from a bare energy bound.  It is
weaker than an \(L^1\) name: the interval depends on the challenge \(g\), and
the construction does not supply one index after which all later levels are
close.

\begin{proposition}[A uniform rate cannot be extracted from the energy bound]
\label{m:no_rate}
There is no computable operation that takes an arbitrary coherent
primitive recursive dyadic mass producer, together with a fixed bound
\(H_n\le5/4\), and always returns an \(L^1\) Cauchy rate for its density steps.
\end{proposition}

\begin{proof}
Fix a machine index \(e\).  At level \(n\), simulate the machine for \(n\)
steps.  If it has not halted at a time \(t<n\), put \(d_n=1\).
If it first halted at \(t<n\), put
\[
 d_n=1+\tfrac12 r_{t+1},
 \]
where \(r_k\) is the level-\(k\) dyadic Rademacher step taking alternating
values \(1,-1\) on sibling cells.  This is a fixed bounded-simulation
algorithm, hence primitive recursive in \(e,n,i\).
At the first nonconstant level, the sibling perturbations average to zero;
all subsequent levels merely refine the same step.  The masses
\(v_{n,i}=2^{-n}d_{n,i}\) are therefore nonnegative, coherent, and normalized.
The energy is either \(1\) or \(5/4\).

If a Cauchy rate returned an index \(N\) valid for distance tolerance \(1/4\),
the machine could not first halt at a time \(t\ge N\): in that case the
distance between \(d_N=1\) and \(d_{t+1}\) would be \(1/2\).
Simulating the machine for \(N\) steps would therefore decide whether it
ever halts.  This contradicts the undecidability of the halting problem.
The argument is a metatheoretic obstruction to a uniform algorithm, not an
unproved claim that one particular numerical sentence is independent of
iPRA.
\end{proof}

\subsection{Object-language proof obligations and efficiency}

The new mathematical proofs reduce to a small number of reusable finite
identities.  For the clipping kernel, the loop invariant after \(j\) cells is
\[
 \sum_{i<j}w_i|a_i-b_i|
 \le 2K\sum_{i<j}w_i\rho(a_i,b_i)
      +\sum_{i<j}w_i t_K(a_i)+\sum_{i<j}w_i t_K(b_i),
 \qquad w_i=2^{-r}.
 \tag{M25}
\]
After denominator clearing this is quantifier-free.  Its base case is zero;
its step adds the proved one-cell inequality.  For refinement, the invariant
says that a processed block of repeated entries contributes the original
entry times its exact multiplicity.  For Fubini, it is equality of two
processed rectangular sums.  For the energy identity, it is the accumulated
identity
\(\sum(d_m-d_n)^2=\sum d_m^2-\sum d_n^2\), together with the
coarse-cell sum identity that cancels the cross term.

The metastability search uses the quantifier-free prefix invariant
\[
 H_{t_j}-H_0=\sum_{i<j}(H_{t_{i+1}}-H_{t_i}),
 \]
and a Boolean accumulator recording whether a good interval has been
found.  All searches have the stated finite bounds.  Proofs of the
unbounded universal conclusions instantiate these finite kernels at
arbitrary precision and level parameters; they do not perform induction on
a formula asserting an infinite convergence statement.

The principal efficiencies are mathematical.  Clipping uses continuous
tests, so sharp-event separation certificates disappear.  The \(L^1\) name
of a dominating function already supplies the required tail data.
Conditional-average convergence obtains its rate from a finite step level.
Radon--Nikodym construction reuses one exact energy identity, and local
sibling bounds discharge the otherwise missing convergence obligation.
The supplied Lean and implementation ledgers separately record which
finite scalar lemmas have been kernel checked and which completion and
object-iPRA derivations remain formalization targets.
